% network.tex - Network and relationship diagram macros for Paperforge
% Requires: _colors.tex, _styles.tex loaded first
%
% Macros provided:
%   \networktriad{center}{node1}{node2}{node3}
%   \networkstar{center}{n1}{n2}{n3}{n4}
%   \dagdiagram - manual DAG construction
%   \venndiagram{set1}{set2}{overlap}
%   \venntriangle{s1}{s2}{s3}{center}
%
% Design: Primary nodes larger, secondary nodes smaller

% ============================================================================
% NETWORK TRIAD (center + 3 connected nodes)
% ============================================================================
% Usage: \networktriad{Center}{Node 1}{Node 2}{Node 3}
\newcommand{\networktriad}[4]{%
\begin{center}
\begin{tikzpicture}
    % Center node (primary)
    \node[diagram node primary] (center) at (0,0) {#1};

    % Surrounding nodes (secondary)
    \node[diagram node secondary] (n1) at (90:2.5cm) {#2};
    \node[diagram node secondary] (n2) at (210:2.5cm) {#3};
    \node[diagram node secondary] (n3) at (330:2.5cm) {#4};

    % Connections
    \draw[diagram arrow] (center) -- (n1);
    \draw[diagram arrow] (center) -- (n2);
    \draw[diagram arrow] (center) -- (n3);
\end{tikzpicture}
\end{center}
}

% ============================================================================
% NETWORK STAR (center + 4 connected nodes)
% ============================================================================
% Usage: \networkstar{Center}{North}{East}{South}{West}
\newcommand{\networkstar}[5]{%
\begin{center}
\begin{tikzpicture}
    % Center node (primary)
    \node[diagram node primary] (center) at (0,0) {#1};

    % Surrounding nodes (secondary)
    \node[diagram node secondary] (n) at (90:2.5cm) {#2};
    \node[diagram node secondary] (e) at (0:2.5cm) {#3};
    \node[diagram node secondary] (s) at (270:2.5cm) {#4};
    \node[diagram node secondary] (w) at (180:2.5cm) {#5};

    % Connections (bidirectional)
    \draw[diagram arrow double] (center) -- (n);
    \draw[diagram arrow double] (center) -- (e);
    \draw[diagram arrow double] (center) -- (s);
    \draw[diagram arrow double] (center) -- (w);
\end{tikzpicture}
\end{center}
}

% ============================================================================
% LINEAR NETWORK (chain of nodes)
% ============================================================================
% Usage: \networkchain{Node1}{Node2}{Node3}{Node4}
\newcommand{\networkchain}[4]{%
\begin{center}
\begin{tikzpicture}[node distance=2.5cm]
    \node[diagram node primary] (n1) {#1};
    \node[diagram node secondary, right=of n1] (n2) {#2};
    \node[diagram node secondary, right=of n2] (n3) {#3};
    \node[diagram node secondary, right=of n3] (n4) {#4};

    \draw[diagram arrow] (n1) -- (n2);
    \draw[diagram arrow] (n2) -- (n3);
    \draw[diagram arrow] (n3) -- (n4);
\end{tikzpicture}
\end{center}
}

% ============================================================================
% FULLY CONNECTED TRIAD
% ============================================================================
% Usage: \networkfully{Node1}{Node2}{Node3}
\newcommand{\networkfully}[3]{%
\begin{center}
\begin{tikzpicture}
    \node[diagram node primary] (n1) at (90:2cm) {#1};
    \node[diagram node primary] (n2) at (210:2cm) {#2};
    \node[diagram node primary] (n3) at (330:2cm) {#3};

    % All connections (bidirectional)
    \draw[diagram arrow double] (n1) -- (n2);
    \draw[diagram arrow double] (n2) -- (n3);
    \draw[diagram arrow double] (n3) -- (n1);
\end{tikzpicture}
\end{center}
}

% ============================================================================
% DIRECTED ACYCLIC GRAPH (3 levels)
% ============================================================================
% Usage: \dagthree{Root}{Mid1}{Mid2}{Leaf1}{Leaf2}{Leaf3}
\newcommand{\dagthree}[6]{%
\begin{center}
\begin{tikzpicture}[node distance=1.5cm]
    % Root level
    \node[diagram node primary] (root) at (0, 0) {#1};

    % Middle level
    \node[diagram node secondary] (m1) at (-2, -1.8) {#2};
    \node[diagram node secondary] (m2) at (2, -1.8) {#3};

    % Leaf level
    \node[diagram node tertiary] (l1) at (-3, -3.6) {#4};
    \node[diagram node tertiary] (l2) at (0, -3.6) {#5};
    \node[diagram node tertiary] (l3) at (3, -3.6) {#6};

    % Edges
    \draw[diagram arrow] (root) -- (m1);
    \draw[diagram arrow] (root) -- (m2);
    \draw[diagram arrow] (m1) -- (l1);
    \draw[diagram arrow] (m1) -- (l2);
    \draw[diagram arrow] (m2) -- (l2);
    \draw[diagram arrow] (m2) -- (l3);
\end{tikzpicture}
\end{center}
}

% ============================================================================
% TWO-SET VENN DIAGRAM
% ============================================================================
% Usage: \venndiagram{Set A}{Set B}{Overlap Label}
\newcommand{\venndiagram}[3]{%
\begin{center}
\begin{tikzpicture}
    % Set A (left circle)
    \draw[draw=diagramPrimary, fill=diagramPrimary!20, thick] (-1,0) circle (1.8cm);
    \node at (-2.2, 0) {\textbf{#1}};

    % Set B (right circle)
    \draw[draw=diagramSecondary, fill=diagramSecondary!20, thick] (1,0) circle (1.8cm);
    \node at (2.2, 0) {\textbf{#2}};

    % Overlap label
    \node at (0, 0) {\small #3};
\end{tikzpicture}
\end{center}
}

% ============================================================================
% THREE-SET VENN DIAGRAM
% ============================================================================
% Usage: \venntriangle{Set A}{Set B}{Set C}{Center Label}
\newcommand{\venntriangle}[4]{%
\begin{center}
\begin{tikzpicture}
    % Set A (top)
    \draw[draw=diagramPrimary, fill=diagramPrimary!15, thick] (0, 1) circle (1.5cm);
    \node at (0, 2.8) {\textbf{#1}};

    % Set B (bottom left)
    \draw[draw=diagramSecondary, fill=diagramSecondary!15, thick] (-1.2, -0.8) circle (1.5cm);
    \node at (-2.5, -1.5) {\textbf{#2}};

    % Set C (bottom right)
    \draw[draw=diagramAccent, fill=diagramAccent!15, thick] (1.2, -0.8) circle (1.5cm);
    \node at (2.5, -1.5) {\textbf{#3}};

    % Center label
    \node at (0, -0.2) {\small\bfseries #4};
\end{tikzpicture}
\end{center}
}

% ============================================================================
% HUB AND SPOKE (business model diagram)
% ============================================================================
% Usage: \hubspoke{Hub}{Spoke1}{Spoke2}{Spoke3}{Spoke4}{Spoke5}{Spoke6}
\newcommand{\hubspoke}[7]{%
\begin{center}
\begin{tikzpicture}
    % Hub (center)
    \node[diagram node primary, minimum size=1.5cm] (hub) at (0,0) {#1};

    % Spokes
    \node[diagram node secondary] (s1) at (60:2.5cm) {#2};
    \node[diagram node secondary] (s2) at (120:2.5cm) {#3};
    \node[diagram node secondary] (s3) at (180:2.5cm) {#4};
    \node[diagram node secondary] (s4) at (240:2.5cm) {#5};
    \node[diagram node secondary] (s5) at (300:2.5cm) {#6};
    \node[diagram node secondary] (s6) at (0:2.5cm) {#7};

    % Connections
    \draw[diagram connector] (hub) -- (s1);
    \draw[diagram connector] (hub) -- (s2);
    \draw[diagram connector] (hub) -- (s3);
    \draw[diagram connector] (hub) -- (s4);
    \draw[diagram connector] (hub) -- (s5);
    \draw[diagram connector] (hub) -- (s6);
\end{tikzpicture}
\end{center}
}

% ============================================================================
% DARK THEME NETWORK
% ============================================================================
% Usage: \networkdark{Center}{Node1}{Node2}{Node3}
\newcommand{\networkdark}[4]{%
\begin{center}
\begin{tikzpicture}
    % Center node
    \node[diagram node dark, fill=diagramDarkPrimary!30] (center) at (0,0) {#1};

    % Surrounding nodes
    \node[diagram node dark] (n1) at (90:2.5cm) {#2};
    \node[diagram node dark] (n2) at (210:2.5cm) {#3};
    \node[diagram node dark] (n3) at (330:2.5cm) {#4};

    % Connections
    \draw[diagram arrow dark] (center) -- (n1);
    \draw[diagram arrow dark] (center) -- (n2);
    \draw[diagram arrow dark] (center) -- (n3);
\end{tikzpicture}
\end{center}
}
